\documentclass{article}
\usepackage{amsmath, amssymb, amsthm}
\usepackage[ruled,vlined]{algorithm2e}
\usepackage[margin=1in]{geometry}

\newtheorem{definition}{Definition}
\newtheorem{lemma}{Lemma}
\newtheorem{theorem}{Theorem}
\newtheorem{proposition}{Proposition}
\newtheorem{remark}{Remark}

\begin{document}

\title{Hybrid Nonlinear Conjugate Gradient and Negative Curvature Descent \\ for Strict Saddle Functions: Complexity Analysis}
\author{}
\date{}
\maketitle


% -------------------- SUMMARY OF WHERE WE ARE, TO UPDATE AS I WORK ON NEW THINGS AND GET NEW RESULTS --------------------

\section*{For Clément and I to see where we are, and where we are going}
In this work, we revisit the Hybrid Gradient and Negative Curvature Descent from~\cite{WrightMa2022}
using the Non-linear Conjugate Gradient with modified restart from~\cite{RoyerChan2022}. 
The first part is an analysis of the complexity of the Hybrid Non-Linear Conjugate Gradient (NCG) and Negative Curvature Descent (NCD) algorithm
for an arbitrary non-convex function. We study this setting under the hypothesis that we know the Lipschitz constants. Therefore, no line search is performed in this case.
This is the trivial case where the complexity of the algorithm is simply the sum of the complexities found in the two papers mentioned:
\begin{equation}
    K = O(\epsilon_g^{-2} + \epsilon_g^{-\max\{1+p, 2(1+p-q)\}} + \epsilon_H^{-3})
\end{equation}

Subsequently, we analyze the hybrid framework under the hypothesis that the non-convex function follows the strict saddle property. 
This setting gives us three regimes to analyze based on the strit saddle conditions:
\begin{itemize}
    \item \textbf{Regime 1}: High gradient. In this case, we apply the standard NCG algorithm considered in~\cite{RoyerChan2022}. We get the same complexity (up to constants):
    \begin{equation}
        K = O(\alpha^{-2} + \alpha^{-\max\{1+p, 2(1 + p - q)\}})
    \end{equation}
    \item \textbf{Regime 2}: Negative Curvature. In this case, we apply the standard NCD algroithm considered in~\cite{WrightMa2022}. We get the same complexity:
    \begin{equation}
        K = O(\beta^{-3})
    \end{equation}
    \item \textbf{Regime 3}: Local strong convexity. This case is where we innovate. It further partitions the analysis into 3 cases:
        \begin{itemize}
        \item The power $1+p-q = 1$: This case gives us the same rate as the restarted case (i.e. Gradient Descent)
        \item The power $1+p-q$ is different from one. This gives us new complexities shown in Section~\ref{sec:strict_saddle}.
    \end{itemize}
\end{itemize}

I was also able to prove that once we enter Regime~3, we stay on the strong convexity region ($B(x^\star, 2\delta)$) thereafter. I actually proved it twice with 2 different techniques
(while trying to prove that entering R3 meant staying in R3), and each lead to condition on different parameters. The first proof, using descent lemma and strong convexity
equations lead to a condition on the Lipschitz constant, with respect to the strog convexity constant. This is not very god since those are hard to control.
The second proof leads to a condition on the $\delta$ parameter, with respect to $\sigma, \kappa, p, q$ and the smoothness constant, which is a bit better. 

However, I was not able to prove that once we enter regime 3,
we stay on regime 3 ($B(x^\star, \delta)$) thereafter, and I think it is impossible to prove (at least starting at iteration $k_0$, the first iteration at which we enter Regime~3).
However, what I am currently working on 

% -------------------- SECTION 1 - SETTING + ASSUMPTIONS --------------------

\section{Setting and Assumptions}
\label{sec:setting}

We consider the unconstrained optimization problem
\begin{equation}
\label{eq:problem}
\min_{x \in \mathbb{R}^n} f(x),
\end{equation}
where $f : \mathbb{R}^n \to \mathbb{R}$ is twice continuously differentiable. We write $g_k = \nabla f(x_k)$ for the gradient at iterate $x_k$.

We make the following standing assumptions throughout.

\begin{description}
    \item[A1.] $\nabla f$ is $L_1$-Lipschitz continuous: $\|\nabla f(x) - \nabla f(y)\| \leq L_1 \|x - y\|$ for all $x, y \in \mathbb{R}^n$.
    \item[A2.] $\nabla^2 f$ is $L_2$-Lipschitz continuous: $\|\nabla^2 f(x) - \nabla^2 f(y)\| \leq L_2 \|x - y\|$ for all $x, y \in \mathbb{R}^n$.
    \item[A3.] $f$ is bounded below: there exists $f_{\text{low}} \in \mathbb{R}$ such that $f(x) \geq f_{\text{low}}$ for all $x \in \mathbb{R}^n$.
\end{description}

We denote the initial gap to the minimum by $\Delta f := f(x_0) - f_{\text{low}}$.

\subsection{Partition of iterations}

Following~\cite{RoyerChan2022}, we partition the iterations into two disjoint sets based on the restart condition. Given parameters $\sigma \in (0,1]$, $\kappa \geq 1$, $p \geq 0$, $q \geq 0$, we define:
\begin{equation}
\label{eq:partition}
\mathcal{N} = \{k \in \mathbb{N} \mid g_k^T d_k \leq -\sigma\|g_k\|^{1+p} \text{ and } \|d_k\| \leq \kappa\|g_k\|^q\}, \qquad \mathcal{R} = \mathbb{N} \setminus \mathcal{N}.
\end{equation}
We call $k \in \mathcal{N}$ a \emph{non-restarted} iteration and $k \in \mathcal{R}$ a \emph{restarted} iteration (where $d_k$ is reset to $-g_k$).

\subsection{Main inequalities}

The following inequalities are used repeatedly. From A1 (descent lemma), for any $x$ and direction $d$:
\begin{equation}
\label{eq:descent_lemma}
f(x + \alpha d) \leq f(x) + \alpha \nabla f(x)^T d + \frac{L_1}{2}\alpha^2\|d\|^2.
\end{equation}
From $L_1$-smoothness, for any $x$ in a neighborhood of a local minimizer $x^\star$:
\begin{equation}
\label{eq:smooth_upper}
\|\nabla f(x)\|^2 \leq 2L_1(f(x) - f(x^\star)).
\end{equation}
If additionally $f$ is $\gamma$-strongly convex in that neighborhood (Polyak--\L{}ojasiewicz inequality):
\begin{equation}
\label{eq:PL}
\|\nabla f(x)\|^2 \geq 2\gamma(f(x) - f(x^\star)).
\end{equation}

% -------------------- SECTION 2 - ALGO --------------------

\section{Algorithm}
\label{sec:algorithm}

The restart condition is:
\begin{equation}
\label{eq:restart}
g_{k+1}^T d_{k+1} \geq -\sigma\|g_{k+1}\|^{1+p} \quad \text{or} \quad \|d_{k+1}\| \geq \kappa\|g_{k+1}\|^q.
\end{equation}

\begin{algorithm}[H]
\caption{Hybrid NCG + Negative Curvature Descent}
\label{algo:hybrid}
\KwIn{$x_0 \in \mathbb{R}^n$, $\sigma \in (0,1]$, $\kappa \geq 1$, $p \geq 0$, $q \geq 0$, $\varepsilon_g > 0$, $\varepsilon_H > 0$.}

\BlankLine
Set $g_0 = \nabla f(x_0)$, $d_0 = -g_0$.

\BlankLine
\For{$k = 0, 1, 2, \ldots$}{

    Compute $g_k = \nabla f(x_k)$.

    \BlankLine
    \eIf{$\|g_k\| \geq \varepsilon_g$}{
        \tcc{NCG phase}
        Compute stepsize $\alpha_k > 0$ (via line search or optimal fixed step).
        
        Set $x_{k+1} = x_k + \alpha_k d_k$ and $g_{k+1} = \nabla f(x_{k+1})$.

        Choose conjugate direction parameter $\beta_{k+1}$.

        Set $d_{k+1} = -g_{k+1} + \beta_{k+1} d_k$.

        \If{condition~\eqref{eq:restart} holds}{
            Restart: $d_{k+1} \leftarrow -g_{k+1}$.
        }
    }{
        Compute smallest eigenvalue $\lambda_k$ and eigenvector $e_k$ of $\nabla^2 f(x_k)$, with sign chosen such that $\langle g_k, e_k \rangle \leq 0$.

        \BlankLine
        \eIf{$-\lambda_k \geq \varepsilon_H$}{
            \tcc{Negative curvature phase}
            $x_{k+1} = x_k + \frac{2\lambda_k}{L_2}\, e_k$.
            
            Reset $d_{k+1} = -\nabla f(x_{k+1})$.
        }{
            \textbf{Stop} and return $x_\star = x_k$.
        }
    }
}

\BlankLine
\textbf{Convergence guarantee:} $\|\nabla f(x_\star)\| \leq \varepsilon_g$, $\;-\lambda_{\min}(\nabla^2 f(x_\star)) \leq \varepsilon_H$.
\end{algorithm}

\paragraph{Stepsize.} Without line search, using the descent lemma~\eqref{eq:descent_lemma} and optimizing over $\alpha$, the optimal fixed stepsize for non-restarted iterations is:
\begin{equation}
\label{eq:stepsize}
\alpha_k^\star = \frac{\sigma}{L_1\kappa^2}\|g_k\|^{1+p-2q},
\end{equation}
while for restarted iterations ($d_k = -g_k$), the stepsize is $\alpha_k = 1/L_1$. Both choices require knowledge of $L_1$ but not a line search. The complexity constants are then:
\begin{equation}
\label{eq:constants}
c_{\mathcal{R}} = \frac{1}{2L_1}, \qquad c_{\mathcal{N}} = \frac{\sigma^2}{2L_1\kappa^2}.
\end{equation}

% -------------------- SECTION 3 - DECREASE LEMMAS --------------------

\section{Per-Iteration Decrease Guarantees}
\label{sec:decrease}

We establish decrease bounds for each type of iteration.

\subsection{NCG phase: non-restarted iterations ($k \in \mathcal{N}$)}

\begin{lemma}
\label{lem:decrease_N}
Let A1 hold and $k \in \mathcal{N}$ with $\|g_k\| > 0$. Then:
\begin{equation}
\label{eq:decrease_N}
f(x_k) - f(x_{k+1}) \geq c_{\mathcal{N}}\, \|g_k\|^{2(1+p-q)}.
\end{equation}
\end{lemma}
\begin{proof}
From the descent lemma~\eqref{eq:descent_lemma}, using $g_k^T d_k \leq -\sigma\|g_k\|^{1+p}$ and $\|d_k\| \leq \kappa\|g_k\|^q$:
\[
f(x_{k+1}) \leq f(x_k) - \alpha\sigma\|g_k\|^{1+p} + \frac{L_1}{2}\alpha^2\kappa^2\|g_k\|^{2q}.
\]
Setting the derivative with respect to $\alpha$ to zero gives $\alpha^\star = \frac{\sigma}{L_1\kappa^2}\|g_k\|^{1+p-2q}$. Substituting:
\begin{align*}
f(x_{k+1}) &\leq f(x_k) - \frac{\sigma^2}{L_1\kappa^2}\|g_k\|^{2(1+p-q)} + \frac{\sigma^2}{2L_1\kappa^2}\|g_k\|^{2(1+p-q)} = f(x_k) - \frac{\sigma^2}{2L_1\kappa^2}\|g_k\|^{2(1+p-q)}. \qedhere
\end{align*}
\end{proof}

\subsection{NCG phase: restarted iterations ($k \in \mathcal{R}$)}

\begin{lemma}
\label{lem:decrease_R}
Let A1 hold and $k \in \mathcal{R}$ with $\|g_k\| > 0$. Then:
\begin{equation}
\label{eq:decrease_R}
f(x_k) - f(x_{k+1}) \geq c_{\mathcal{R}}\, \|g_k\|^{2}.
\end{equation}
\end{lemma}
\begin{proof}
For restarted iterations, $d_k = -g_k$, corresponding to $\sigma = \kappa = 1$ and $p = q = 1$ in Lemma~\ref{lem:decrease_N}. The stepsize $\alpha = 1/L_1$ gives:
\[
f(x_{k+1}) \leq f(x_k) - \frac{1}{2L_1}\|g_k\|^2. \qedhere
\]
\end{proof}

\subsection{Negative curvature phase}

\begin{lemma}
\label{lem:decrease_NCD}
Let A1--A2 hold and $-\lambda_k \geq \varepsilon_H > 0$. Then the step $x_{k+1} = x_k + \frac{2\lambda_k}{L_2}e_k$ satisfies:
\begin{equation}
\label{eq:decrease_NCD}
f(x_k) - f(x_{k+1}) \geq \frac{2|\lambda_k|^3}{3L_2^2} \geq \frac{2\varepsilon_H^3}{3L_2^2}.
\end{equation}
\end{lemma}
\begin{proof}
Since $\nabla^2 f$ is $L_2$-Lipschitz, a third-order Taylor expansion gives:
\[
f(x_{k+1}) \leq f(x_k) + \langle g_k, x_{k+1} - x_k \rangle + \frac{1}{2}\langle \nabla^2 f(x_k)(x_{k+1} - x_k), x_{k+1} - x_k \rangle + \frac{L_2}{6}\|x_{k+1} - x_k\|^3.
\]
Since $e_k$ is a unit eigenvector with eigenvalue $\lambda_k$ and $\langle g_k, e_k \rangle \leq 0$, $\lambda_k < 0$, the first-order term $\frac{2\lambda_k}{L_2}\langle g_k, e_k\rangle \geq 0$. Dropping it:
\begin{align*}
f(x_{k+1}) &\leq f(x_k) + \frac{1}{2}\cdot\frac{2\lambda_k^2}{L_2}\cdot\frac{2\lambda_k}{L_2} + \frac{L_2}{6}\left(\frac{2|\lambda_k|}{L_2}\right)^3 = f(x_k) - \frac{2|\lambda_k|^3}{3L_2^2}. \qedhere
\end{align*}
\end{proof}

\begin{remark}
\label{rem:monotone}
All three lemmas guarantee $f(x_{k+1}) < f(x_k)$. This monotone decrease property ensures iterates remain within sublevel sets of $f$.
\end{remark}

% -------------------- SECTION 4 - ARBITRARY NONCONVEX - WE MIGHT DELETE THIS LATER, IT WAS MORE OF AN EXAMPLE AND EXERCISE FOR ME --------------------

\section{Complexity on Arbitrary Nonconvex Functions}
\label{sec:arbitrary}

We first analyze Algorithm~\ref{algo:hybrid} on a general nonconvex function satisfying A1--A3, without assuming the strict saddle property.

\begin{theorem}
\label{thm:arbitrary}
Under A1--A3, Algorithm~\ref{algo:hybrid} reaches a point satisfying $\|\nabla f(x_\star)\| \leq \varepsilon_g$ and $-\lambda_{\min}(\nabla^2 f(x_\star)) \leq \varepsilon_H$ in at most $K$ iterations, where:
\begin{equation}
\label{eq:bound_arbitrary}
K \leq \frac{\Delta f}{c_{\mathcal{R}}}\,\varepsilon_g^{-2} + \frac{\Delta f}{c_{\mathcal{N}}}\,\varepsilon_g^{-2(1+p-q)} + \frac{3L_2^2\,\Delta f}{2}\,\varepsilon_H^{-3}.
\end{equation}
\end{theorem}

\begin{proof}
Let $K$ be the first index such that the stopping condition holds. Every iteration $k < K$ is either an NCG step ($\|g_k\| \geq \varepsilon_g$) 
or an NCD step ($\|g_k\| < \varepsilon_g$, $-\lambda_k \geq \varepsilon_H$). Among NCG iterations, each is either restarted or non-restarted.

By Lemma~\ref{lem:decrease_R}, each restarted NCG iteration decreases $f$ by at least $c_{\mathcal{R}}\,\varepsilon_g^2$. 
By Lemma~\ref{lem:decrease_N}, each non-restarted NCG iteration decreases $f$ by at least $c_{\mathcal{N}}\,\varepsilon_g^{2(1+p-q)}$. By Lemma~\ref{lem:decrease_NCD}, each NCD iteration decreases $f$ by at least $\frac{2\varepsilon_H^3}{3L_2^2}$.

Telescoping and using A3:
\[
\Delta f \geq K^1_{\mathcal{R}} \cdot c_{\mathcal{R}}\,\varepsilon_g^2 + K^1_{\mathcal{N}} \cdot c_{\mathcal{N}}\,\varepsilon_g^{2(1+p-q)} + K^2 \cdot \frac{2\varepsilon_H^3}{3L_2^2}.
\]
Since all terms are positive, each is individually bounded by $\Delta f$. Summing the resulting bounds on $K^1_{\mathcal{R}}$, $K^1_{\mathcal{N}}$, and $K^2$ gives~\eqref{eq:bound_arbitrary}.
\end{proof}

\begin{remark}
The bound~\eqref{eq:bound_arbitrary} is conservative because each group uses the full budget $\Delta f$ independently. The NCG terms match those of Theorem~3.1 in~\cite{RoyerChan2022} (up to constants due to the absence of line search), while the NCD term follows from~\cite{WrightMa2022}.
\end{remark}

% -------------------- SECTION 5 - STRICT SADDLE FUNCTION - THE HEART OF THE ANALYSIS --------------------

\section{Complexity on Strict Saddle Functions}
\label{sec:strict_saddle}

We now specialize to functions satisfying the strict saddle property, which provides additional structure near stationary points.

\subsection{The strict saddle property}

\begin{definition}[\cite{GeSaddle15}]
\label{def:strict_saddle}
A function $f$ is $(\zeta, \beta, \gamma, \delta)$-strict saddle if, for any $x \in \mathbb{R}^n$, at least one of the following holds:
\begin{enumerate}
    \item \textbf{(Large gradient)} $\|\nabla f(x)\| \geq \zeta$,
    \item \textbf{(Negative curvature)} $\lambda_{\min}(\nabla^2 f(x)) \leq -\beta$,
    \item \textbf{(Near a local minimizer)} There exists a local minimizer $x^\star$ with $\|x - x^\star\| \leq \delta$, and $f$ restricted to $\{x' : \|x' - x^\star\| \leq 2\delta\}$ is $\gamma$-strongly convex.
\end{enumerate}
\end{definition}

Setting $\varepsilon_g \leq \zeta$ and $\varepsilon_H \leq \beta$ in Algorithm~\ref{algo:hybrid} ensures that upon termination, condition~(3) holds: the output is near a local minimizer. The three conditions define three regimes, which we analyze separately.

\paragraph{Notation.} We partition the iteration indices $\{0, \ldots, K-1\}$ into five disjoint sets:
\begin{center}
\begin{tabular}{lll}
$\mathcal{N}^1$: & non-restarted in Regime 1 & ($|\mathcal{N}^1| = K^1_{\mathcal{N}}$) \\
$\mathcal{R}^1$: & restarted in Regime 1 & ($|\mathcal{R}^1| = K^1_{\mathcal{R}}$) \\
$\mathcal{H}$: & negative curvature in Regime 2 & ($|\mathcal{H}| = K^2$) \\
$\mathcal{N}^3$: & non-restarted in Regime 3 & ($|\mathcal{N}^3| = K^3_{\mathcal{N}}$) \\
$\mathcal{R}^3$: & restarted in Regime 3 & ($|\mathcal{R}^3| = K^3_{\mathcal{R}}$)
\end{tabular}
\end{center}
The total iteration count is $K = K^1_{\mathcal{N}} + K^1_{\mathcal{R}} + K^2 + K^3_{\mathcal{N}} + K^3_{\mathcal{R}}$.

All iterations draw from the single decrease budget $\Delta f$:
\begin{equation}
\label{eq:budget}
\begin{split}
    \Delta f \geq{} &\sum_{k \in \mathcal{N}^1}(f(x_k) - f(x_{k+1})) + \sum_{k \in \mathcal{R}^1}(f(x_k) - f(x_{k+1})) + \sum_{k \in \mathcal{H}}(f(x_k) - f(x_{k+1})) \\
    &+ \sum_{k \in \mathcal{N}^3}(f(x_k) - f(x_{k+1})) + \sum_{k \in \mathcal{R}^3}(f(x_k) - f(x_{k+1})).
\end{split}
\end{equation}
We define the \emph{decrease consumed} in each regime as the sum of per-iteration decreases over all iterations belonging to that regime:
\[
\Delta f^1 := \sum_{k \in \mathcal{N}^1 \cup \mathcal{R}^1}(f(x_k) - f(x_{k+1})), \qquad \Delta f^2 := \sum_{k \in \mathcal{H}}(f(x_k) - f(x_{k+1})), \qquad \Delta f^3 := \sum_{k \in \mathcal{N}^3 \cup \mathcal{R}^3}(f(x_k) - f(x_{k+1})),
\]
so that $\Delta f^1 + \Delta f^2 + \Delta f^3 \leq \Delta f$. Note that $\Delta f^j$ is not the suboptimality at entry to regime $j$, but the total amount of function decrease accumulated during all iterations spent in regime~$j$. In particular, since the algorithm may visit regimes in any order, the iterations within each regime need not be contiguous.

For Regime~3, however, the monotone decrease property (Remark~\ref{rem:monotone}) together with the analysis in Section~\ref{subsec:regime3} guarantees that once the algorithm enters Regime~3 at some iteration $k_0$, it stays there. In this case, $\Delta f^3 = f(x_{k_0}) - f(x_K)$, where $K$ is the final iteration. Since $f(x_K) \geq f(x^\star)$:
\[
\Delta f^3 \leq f(x_{k_0}) - f(x^\star) \leq \frac{L_1}{2}\delta^2,
\]
where the last inequality follows from~\eqref{eq:entry_gap}. This trajectory-independent upper bound on $\Delta f^3$ is what allows $\delta$ to appear in the Regime~3 complexity bounds.

\subsection{Regime 1: Large gradient ($\|g_k\| \geq \zeta$)}

The per-iteration decrease bounds from Lemmas~\ref{lem:decrease_N} and~\ref{lem:decrease_R} give:
\[
K^1_{\mathcal{R}} \leq \frac{\Delta f^1}{c_{\mathcal{R}}\,\zeta^2}, \qquad K^1_{\mathcal{N}} \leq \frac{\Delta f^1}{c_{\mathcal{N}}\,\zeta^{2(1+p-q)}}.
\]
Since $\Delta f^1 \leq \Delta f$:
\begin{equation}
\label{eq:regime1}
K^1 = K^1_{\mathcal{R}} + K^1_{\mathcal{N}} \leq \frac{\Delta f}{c_{\mathcal{R}}}\,\zeta^{-2} + \frac{\Delta f}{c_{\mathcal{N}}}\,\zeta^{-2(1+p-q)}.
\end{equation}

\subsection{Regime 2: Negative curvature ($\lambda_{\min}(\nabla^2 f(x_k)) \leq -\beta$)}

From Lemma~\ref{lem:decrease_NCD} and $\Delta f^2 \leq \Delta f$:
\begin{equation}
\label{eq:regime2}
K^2 \leq \frac{3L_2^2\,\Delta f}{2\beta^3}.
\end{equation}

\subsection{Regime 3: Near a local minimizer}
\label{subsec:regime3}

This is where $\|x_k - x^\star\| \leq \delta$ and $f$ is $\gamma$-strongly convex in the $2\delta$-neighborhood. Both~\eqref{eq:smooth_upper} and~\eqref{eq:PL} hold. The parameter $\delta$ controls the initial suboptimality upon entering this regime.

\subsubsection{Entry condition and the role of $\delta$}
\label{subsubsec:entry}

At the first iteration $k_0$ in Regime~3, the strict saddle property (Definition~\ref{def:strict_saddle}, condition 3) guarantees $\|x_{k_0} - x^\star\| \leq \delta$ for some local minimizer $x^\star$. We now derive bounds on the suboptimality gap $\Delta f^3 := f(x_{k_0}) - f(x^\star)$ and the gradient norm $\|g_{k_0}\|$ in terms of $\delta$.

\paragraph{Upper bound on the suboptimality gap.} Since $\nabla f$ is $L_1$-Lipschitz (A1), the descent lemma~\eqref{eq:descent_lemma} gives, for any $x, y \in \mathbb{R}^n$:
\[
f(y) \leq f(x) + \nabla f(x)^T(y - x) + \frac{L_1}{2}\|y - x\|^2.
\]
Setting $x = x^\star$ and $y = x_{k_0}$, and using $\nabla f(x^\star) = 0$ (since $x^\star$ is a local minimizer):
\[
f(x_{k_0}) \leq f(x^\star) + 0 + \frac{L_1}{2}\|x_{k_0} - x^\star\|^2 \leq f(x^\star) + \frac{L_1}{2}\delta^2.
\]
Therefore:
\begin{equation}
\label{eq:entry_upper}
\Delta f^3 = f(x_{k_0}) - f(x^\star) \leq \frac{L_1}{2}\delta^2.
\end{equation}

\paragraph{Lower bound on the suboptimality gap.} Since $f$ is $\gamma$-strongly convex in $B(x^\star, 2\delta)$ and $x_{k_0} \in B(x^\star, \delta) \subset B(x^\star, 2\delta)$:
\[
f(x_{k_0}) \geq f(x^\star) + \nabla f(x^\star)^T(x_{k_0} - x^\star) + \frac{\gamma}{2}\|x_{k_0} - x^\star\|^2 = f(x^\star) + \frac{\gamma}{2}\|x_{k_0} - x^\star\|^2,
\]
where we again used $\nabla f(x^\star) = 0$. Therefore:
\begin{equation}
\label{eq:entry_lower}
\Delta f^3 = f(x_{k_0}) - f(x^\star) \geq \frac{\gamma}{2}\|x_{k_0} - x^\star\|^2.
\end{equation}
Combining~\eqref{eq:entry_upper} and~\eqref{eq:entry_lower}:
\begin{equation}
\label{eq:entry_gap}
\frac{\gamma}{2}\|x_{k_0} - x^\star\|^2 \leq \Delta f^3 \leq \frac{L_1}{2}\delta^2.
\end{equation}

\paragraph{Upper bound on the gradient norm.} From the smoothness inequality~\eqref{eq:smooth_upper}, which states $\|\nabla f(x)\|^2 \leq 2L_1(f(x) - f(x^\star))$ 
for any $x$ near $x^\star$, we have:
\[
\|g_{k_0}\|^2 \leq 2L_1(f(x_{k_0}) - f(x^\star)).
\]
Substituting the upper bound~\eqref{eq:entry_upper} on $\Delta f^3$:
\[
\|g_{k_0}\|^2 \leq 2L_1 \cdot \frac{L_1}{2}\delta^2 = L_1^2\delta^2.
\]
Taking square roots:
\begin{equation}
\label{eq:entry_grad_upper}
\|g_{k_0}\| \leq L_1\delta.
\end{equation}

\paragraph{Lower bound on the gradient norm.} From the Polyak--\L{}ojasiewicz inequality~\eqref{eq:PL}, which states \\
 $\|\nabla f(x)\|^2 \geq 2\gamma(f(x) - f(x^\star))$, 
and the lower bound~\eqref{eq:entry_lower}:
\[
\|g_{k_0}\|^2 \geq 2\gamma(f(x_{k_0}) - f(x^\star)) \geq 2\gamma \cdot \frac{\gamma}{2}\|x_{k_0} - x^\star\|^2 = \gamma^2\|x_{k_0} - x^\star\|^2.
\]
Taking square roots:
\begin{equation}
\label{eq:entry_grad_lower}
\|g_{k_0}\| \geq \gamma\|x_{k_0} - x^\star\|.
\end{equation}

\subsubsection{Non-restarted iterations in Regime 3 ($k \in \mathcal{N}^3$)}

\paragraph{Step 1: Combine decrease lemma with PL inequality.} Lemma~\ref{lem:decrease_N} gives:
\[
f(x_k) - f(x_{k+1}) \geq c_{\mathcal{N}}\|g_k\|^{2r}, \qquad r := 1 + p - q.
\]
The PL inequality~\eqref{eq:PL} states $\|g_k\|^2 \geq 2\gamma(f(x_k) - f(x^\star))$. Substituting:
\[
f(x_k) - f(x_{k+1}) \geq c_{\mathcal{N}}\left(2\gamma(f(x_k) - f(x^\star))\right)^r.
\]
Let $C := c_{\mathcal{N}}(2\gamma)^r = \frac{\sigma^2(2\gamma)^r}{2L_1\kappa^2}$. Then:
\begin{equation}
\label{eq:decrease_N3}
f(x_k) - f(x_{k+1}) \geq C(f(x_k) - f(x^\star))^r.
\end{equation}

\paragraph{Step 2: Transform to a bound on $f(x_{k+1}) - f(x^\star)$.} Rearranging~\eqref{eq:decrease_N3}:
\begin{equation}
\label{eq:contraction_N3}
f(x_{k+1}) - f(x^\star) \leq (f(x_k) - f(x^\star)) - C(f(x_k) - f(x^\star))^r = (f(x_k) - f(x^\star))\left(1 - C(f(x_k) - f(x^\star))^{r-1}\right).
\end{equation}

We also need to relate the stopping criterion $\|g_k\| \leq \varepsilon_g$ to a condition on $f(x_k) - f(x^\star)$. The smoothness bound~\eqref{eq:smooth_upper} gives $\|g_k\|^2 \leq 2L_1(f(x_k) - f(x^\star))$, so $\|g_k\| \leq \varepsilon_g$ is guaranteed when $f(x_k) - f(x^\star) \leq \varepsilon_g^2/(2L_1)$. We therefore set:
\begin{equation}
\label{eq:eps_convert}
\varepsilon := \frac{\varepsilon_g^2}{2L_1},
\end{equation}
and the stopping criterion becomes $f(x_k) - f(x^\star) \leq \varepsilon$.

We now analyze three cases depending on $r$.

\paragraph{Step 3: Case 1, $r = 1$ ($p = q$).} The recursion~\eqref{eq:contraction_N3} becomes:
\[
f(x_{k+1}) - f(x^\star) \leq (1 - C)(f(x_k) - f(x^\star)).
\]
Since $C = \frac{\sigma^2(2\gamma)}{2L_1\kappa^2} < 1$, this is a geometric contraction. Iterating $K$ times:
\[
f(x_K) - f(x^\star) \leq (1 - C)^K(f(x_{k_0}) - f(x^\star)).
\]
Requiring $f(x_K) - f(x^\star) \leq \varepsilon$:
\[
(1 - C)^K \leq \frac{\varepsilon}{f(x_{k_0}) - f(x^\star)}.
\]
Taking logarithms. Since $1 - C < 1$, $\log(1 - C) < 0$, so dividing flips the inequality:
\[
K \geq \frac{\log\left(\frac{\varepsilon}{f(x_{k_0}) - f(x^\star)}\right)}{\log(1 - C)} = \frac{\log\left(\frac{f(x_{k_0}) - f(x^\star)}{\varepsilon}\right)}{-\log(1 - C)}.
\]
Using $-\log(1 - C) \geq C$ for $C \in (0,1)$, and $\Delta f^3 = f(x_{k_0}) - f(x^\star) \leq \frac{L_1}{2}\delta^2$:
\[
K \geq \frac{1}{C}\log\frac{\Delta f^3}{\varepsilon}.
\]
Substituting $C = \frac{\sigma^2(2\gamma)}{2L_1\kappa^2}$ and $\varepsilon = \frac{\varepsilon_g^2}{2L_1}$:
\begin{equation}
\label{eq:N3_r1}
K^3_{\mathcal{N},\, p=q} \leq \frac{L_1\kappa^2}{\sigma^2\gamma}\log\frac{L_1^2\delta^2}{\varepsilon_g^2}.
\end{equation}

\paragraph{Step 4: Case 2, $r > 1$ ($p > q$, sublinear rate).} From~\eqref{eq:contraction_N3}:
\[
f(x_{k+1}) - f(x^\star) \leq (f(x_k) - f(x^\star))\left(1 - C(f(x_k) - f(x^\star))^{r-1}\right).
\]
We first verify that $x := C(f(x_k) - f(x^\star))^{r-1} \in (0,1)$. Positivity is clear ($C > 0$ and $f(x_k) > f(x^\star)$). If $x \geq 1$, the right-hand side is $\leq 0$, giving $f(x_{k+1}) \leq f(x^\star)$. But since $x_{k+1} \in B(x^\star, 2\delta)$ where $f$ is strongly convex, $f(x_{k+1}) \geq f(x^\star)$. Therefore $f(x_{k+1}) = f(x^\star)$: the algorithm has already converged. So for a non-terminated run, $x \in (0,1)$.

Define the potential function $\phi_k := (f(x_k) - f(x^\star))^{1-r}$, where $1 - r < 0$. Raising both sides of~\eqref{eq:contraction_N3} to the power $1 - r < 0$ flips the inequality:
\[
\phi_{k+1} \geq (1 - C\phi_k^{-1})^{1-r}\,\phi_k.
\]
Using the inequality $(1 - x)^{1-r} \geq 1 + (r-1)x$ for $x \in (0,1)$:
\begin{align*}
\phi_{k+1} &\geq \phi_k\left(1 + (r-1)C\,\phi_k^{-1}\right) = \phi_k + (r-1)C.
\end{align*}
Telescoping over $K^3_{\mathcal{N}}$ non-restarted iterations:
\[
\phi_K \geq \phi_0 + K^3_{\mathcal{N}}(r-1)C.
\]
Requiring $f(x_K) - f(x^\star) \leq \varepsilon$. Since $1 - r < 0$, this gives $\phi_K = (f(x_K) - f(x^\star))^{1-r} \geq \varepsilon^{1-r}$. So we need:
\[
\phi_0 + K^3_{\mathcal{N}}(r-1)C \geq \varepsilon^{1-r},
\]
hence:
\[
K^3_{\mathcal{N}} \leq \frac{\varepsilon^{1-r} - \phi_0}{(r-1)C}.
\]
Substituting $\phi_0 = (\Delta f^3)^{1-r} \geq (L_1\delta^2/2)^{1-r}$, $C = \frac{\sigma^2(2\gamma)^r}{2L_1\kappa^2}$, and $\varepsilon = \frac{\varepsilon_g^2}{2L_1}$:
\begin{equation}
\label{eq:N3_r_big}
K^3_{\mathcal{N},\, p>q} \leq \frac{2L_1\kappa^2}{\sigma^2(2\gamma)^{1+p-q}(p-q)}\left(\left(\frac{\varepsilon_g^2}{2L_1}\right)^{q-p} - \left(\frac{L_1\delta^2}{2}\right)^{q-p}\right).
\end{equation}

\paragraph{Step 5: Case 3, $r < 1$ ($p < q$, finite termination).} From~\eqref{eq:contraction_N3}:
\[
f(x_{k+1}) - f(x^\star) \leq (f(x_k) - f(x^\star))\left(1 - C(f(x_k) - f(x^\star))^{r-1}\right).
\]
Define the potential function $\phi_k := (f(x_k) - f(x^\star))^{1-r}$, where $1 - r > 0$. Raising both sides of~\eqref{eq:contraction_N3} to the power $1 - r > 0$ preserves the inequality:
\[
\phi_{k+1} \leq (1 - C\phi_k^{-1})^{1-r}\,\phi_k.
\]
Using the inequality $(1 - x)^{1-r} \geq 1 + (r-1)x$ for $x \in (0,1)$:
\begin{align*}
\phi_{k+1} &\leq \phi_k\left(1 + (r-1)C\,\phi_k^{-1}\right) = \phi_k + (r-1)C.
\end{align*}
Since $r - 1 < 0$, this reads $\phi_{k+1} \leq \phi_k - (1-r)C$, i.e., $\phi_k$ decreases by a fixed amount at each iteration.

Telescoping over $K^3_{\mathcal{N}}$ non-restarted iterations:
\[
\phi_K \leq \phi_0 - K^3_{\mathcal{N}}(1-r)C.
\]
Since $\phi_K = (f(x_K) - f(x^\star))^{1-r} \geq 0$ (a nonnegative quantity raised to a positive power):
\[
0 \leq \phi_0 - K^3_{\mathcal{N}}(1-r)C \implies K^3_{\mathcal{N}} \leq \frac{\phi_0}{(1-r)C}.
\]
At $K^\star = \lfloor \frac{\phi_0}{(1-r)C} \rfloor$, the recursion requires $\phi_{K^\star} \leq 0$. Since $\phi_{K^\star} \geq 0$, the only possibility is $\phi_{K^\star} = 0$, which implies $f(x_{K^\star}) = f(x^\star)$ exactly.

Substituting $\phi_0 = (\Delta f^3)^{1-r} \leq (L_1\delta^2/2)^{1-r}$ and $C = \frac{\sigma^2(2\gamma)^r}{2L_1\kappa^2}$:
\begin{equation}
\label{eq:N3_r_small}
K^3_{\mathcal{N},\, p<q} \leq \left(\frac{L_1\delta^2}{2}\right)^{q-p}\frac{2L_1\kappa^2}{\sigma^2(2\gamma)^{1+p-q}(q-p)}.
\end{equation}
This bound is independent of $\varepsilon_g$: the algorithm reaches the exact minimizer in finitely many non-restarted steps.

\subsubsection{Restarted iterations in Regime 3 ($k \in \mathcal{R}^3$)}

From Lemma~\ref{lem:decrease_R} and the PL inequality:
\[
f(x_{k+1}) - f(x^\star) \leq \left(1 - \frac{\gamma}{L_1}\right)(f(x_k) - f(x^\star)).
\]
Iterating and requiring $f(x_K) - f(x^\star) \leq \varepsilon$, using $-\log(1-x) \geq x$ for $x \in (0,1)$:
\begin{equation}
\label{eq:R3}
K^3_{\mathcal{R}} \leq \frac{L_1}{\gamma}\log\frac{L_1\delta^2}{2\varepsilon}.
\end{equation}

\subsection{Overall complexity}

\begin{theorem}
\label{thm:strict_saddle}
Under A1--A3, suppose $f$ is $(\zeta, \beta, \gamma, \delta)$-strict saddle. Then Algorithm~\ref{algo:hybrid} with $\varepsilon_g = \zeta$ and $\varepsilon_H = \beta$ reaches a point near a local minimizer in at most $K$ iterations, where:
\begin{equation}
\label{eq:total}
\boxed{
K \leq \underbrace{\frac{\Delta f}{c_{\mathcal{R}}}\,\zeta^{-2} + \frac{\Delta f}{c_{\mathcal{N}}}\,\zeta^{-2(1+p-q)}}_{\text{Regime 1}} + \underbrace{\frac{3L_2^2\,\Delta f}{2\beta^3}}_{\text{Regime 2}} + \underbrace{\frac{L_1}{\gamma}\log\frac{L_1\delta^2}{2\varepsilon} + K^3_{\mathcal{N}}(r)}_{\text{Regime 3}}
}
\end{equation}
where $K^3_{\mathcal{N}}(r)$ depends on $r = 1+p-q$:
\[
K^3_{\mathcal{N}}(r) = 
\begin{cases}
\dfrac{L_1\kappa^2}{\sigma^2\gamma}\log\dfrac{L_1\delta^2}{2\varepsilon} & \text{if } r = 1, \\[12pt]
\dfrac{2L_1\kappa^2}{\sigma^2(2\gamma)^r(r-1)}\left(\varepsilon^{1-r} - \left(\dfrac{L_1\delta^2}{2}\right)^{1-r}\right) & \text{if } r > 1, \\[12pt]
\dfrac{2L_1\kappa^2}{\sigma^2(2\gamma)^r(1-r)}\left(\dfrac{L_1\delta^2}{2}\right)^{1-r} & \text{if } r < 1.
\end{cases}
\]
\end{theorem}

\begin{proof}
Combine~\eqref{eq:regime1},~\eqref{eq:regime2}, and the Regime~3 bounds from Sections~\ref{subsec:regime3}. The bound is conservative because each regime 
uses $\Delta f^j \leq \Delta f$ independently, whereas $\Delta f^1 + \Delta f^2 + \Delta f^3 \leq \Delta f$. If we knew all $\Delta f^j$ in advance, we could
plug it in the bound to have something tighter.
\end{proof}

The restarted terms ($\zeta^{-2}$ in Regime~1, $\frac{L_1}{\gamma}\log\frac{1}{\varepsilon}$ in Regime~3) match GD. 
The overall bound strictly improves over GD only when $K^1_{\mathcal{R}} + K^3_{\mathcal{R}} \ll K$. 

% ------------------- Analysis of staying in REGIME 3 (ONGOING) -------------------

\section{Regime 3 Analysis}

\subsection{Staying in Strong-Convexity Region (proof 1 using strongly convex + descent lemma)}
\label{subsubsec:staying}

\begin{proposition}
\label{prop:stay}
Suppose $f$ is $\gamma$-strongly convex in $B(x^\star, 2\delta) := \{x : \|x - x^\star\| \leq 2\delta\}$ and A1 holds globally. Let $k_0$ be the first iteration at which $x_{k_0} \in B(x^\star, \delta)$. If $L_1 \leq 4\gamma$, then $x_k \in B(x^\star, 2\delta)$ for all $k \geq k_0$.
\end{proposition}

\begin{proof}
We proceed by induction on $k \geq k_0$.

\paragraph{Base case ($k = k_0$).} By assumption, $x_{k_0} \in B(x^\star, \delta)$. Since $\delta < 2\delta$, we have $x_{k_0} \in B(x^\star, 2\delta)$.

\paragraph{Inductive step.} Assume $x_k \in B(x^\star, 2\delta)$ for some $k \geq k_0$. We want to show $x_{k+1} \in B(x^\star, 2\delta)$.

\medskip
\emph{Step 1: $f(x_{k+1}) \leq f(x_{k_0})$.} Since A1 holds globally, the descent lemma~\eqref{eq:descent_lemma} applies at every iterate, regardless of its location. As established in Lemmas~\ref{lem:decrease_N},~\ref{lem:decrease_R}, and~\ref{lem:decrease_NCD}, every iteration of Algorithm~\ref{algo:hybrid} (whether non-restarted, restarted, or negative curvature) produces strict decrease:
\[
f(x_{k+1}) < f(x_k) \qquad \text{for all } k \geq 0.
\]
Applying this repeatedly from $k_0$ to $k$:
\[
f(x_{k+1}) < f(x_k) < f(x_{k-1}) < \cdots < f(x_{k_0}).
\]
In particular:
\begin{equation}
\label{eq:chain}
f(x_{k+1}) \leq f(x_{k_0}).
\end{equation}

\medskip
\emph{Step 2: Upper bound on $f(x_{k_0})$.} From the entry condition (Section~\ref{subsubsec:entry}), since $x_{k_0} \in B(x^\star, \delta)$ and $\nabla f(x^\star) = 0$, the descent lemma with $x = x^\star$ and $y = x_{k_0}$ gives:
\[
f(x_{k_0}) \leq f(x^\star) + \frac{L_1}{2}\|x_{k_0} - x^\star\|^2 \leq f(x^\star) + \frac{L_1}{2}\delta^2.
\]
Combining with~\eqref{eq:chain}:
\begin{equation}
\label{eq:fk1_bound}
f(x_{k+1}) \leq f(x^\star) + \frac{L_1}{2}\delta^2.
\end{equation}

\medskip
\emph{Step 3: Lower bound on $f$ at the boundary of $B(x^\star, 2\delta)$.} Suppose for contradiction that $\|x_{k+1} - x^\star\| > 2\delta$, i.e., $x_{k+1} \notin B(x^\star, 2\delta)$. Consider any point $\bar{x}$ on the boundary of $B(x^\star, 2\delta)$, i.e., $\|\bar{x} - x^\star\| = 2\delta$. Since $f$ is $\gamma$-strongly convex in $B(x^\star, 2\delta)$ and $\nabla f(x^\star) = 0$:
\[
f(\bar{x}) \geq f(x^\star) + \nabla f(x^\star)^T(\bar{x} - x^\star) + \frac{\gamma}{2}\|\bar{x} - x^\star\|^2 = f(x^\star) + \frac{\gamma}{2}(2\delta)^2 = f(x^\star) + 2\gamma\delta^2.
\]
This means: to leave $B(x^\star, 2\delta)$, the function value must exceed $f(x^\star) + 2\gamma\delta^2$, since $f$ increases (by strong convexity) as we move away from $x^\star$. By continuity of $f$, if $x_{k+1}$ is outside $B(x^\star, 2\delta)$, there exists a point on the path from $x_k$ to $x_{k+1}$ where $f$ equals $f(x^\star) + 2\gamma\delta^2$. Since $f(x_{k+1}) < f(x_k)$ and $x_k \in B(x^\star, 2\delta)$ (inductive hypothesis), the iterate must cross the boundary, so:
\begin{equation}
\label{eq:boundary}
f(x_{k+1}) \geq f(x^\star) + 2\gamma\delta^2.
\end{equation}

\medskip
\emph{Step 4: Contradiction.} Combining~\eqref{eq:fk1_bound} and~\eqref{eq:boundary}:
\[
f(x^\star) + 2\gamma\delta^2 \leq f(x_{k+1}) \leq f(x^\star) + \frac{L_1}{2}\delta^2.
\]
Subtracting $f(x^\star)$ from all terms:
\[
2\gamma\delta^2 \leq \frac{L_1}{2}\delta^2.
\]
Dividing by $\delta^2 > 0$:
\[
2\gamma \leq \frac{L_1}{2} \iff 4\gamma \leq L_1.
\]
This contradicts the assumption $L_1 \leq 4\gamma$. Therefore $x_{k+1} \in B(x^\star, 2\delta)$, completing the inductive step.

\end{proof}


\subsection{Staying in Strong-Convexity Region (proof 2 using step length bound)}

We provide an alternative proof that iterates remain in $B(x^\star, 2\delta)$, using control of the step length. 
This approach replaces the condition $L_1 \leq 4\gamma$ with a condition on $\delta$.

\begin{proposition}
\label{prop:stay_step}
Suppose $f$ is $\gamma$-strongly convex in $B(x^\star, 2\delta)$ and A1 holds globally. Let $k_0$ be the first iteration at which $x_{k_0} \in B(x^\star, \delta)$. Suppose one of the following holds:
\begin{enumerate}
    \item $p = q$,
    \item $p > q$ and $\delta \leq \left(\dfrac{\kappa}{\sigma}\right)^{1/(p-q)}\dfrac{1}{L_1}$,
    \item $p < q$ and $\delta \geq \left(\dfrac{\sigma}{\kappa}\right)^{1/(q-p)}\dfrac{1}{L_1}$.
\end{enumerate}
Then $x_k \in B(x^\star, 2\delta)$ for all $k \geq k_0$.
\end{proposition}

\begin{proof}
We proceed by induction on $k \geq k_0$. The base case $x_{k_0} \in B(x^\star, \delta) \subset B(x^\star, 2\delta)$ holds by assumption. Assume $x_k \in B(x^\star, 2\delta)$ for some $k \geq k_0$. We show $x_{k+1} \in B(x^\star, 2\delta)$ by bounding the step length for each iteration type.

\subsubsection{Non-restarted case ($k \in \mathcal{N}$)}

Recall that $x_{k+1} = x_k + \alpha_k d_k$, where $d_k = -g_k + \beta_k d_{k-1}$ (with $d_0 = -g_0$). Under non-restarted iterations, condition~\eqref{eq:restart} is not triggered, so:
\[
g_{k+1}^T d_{k+1} \leq -\sigma\|g_{k+1}\|^{1+p} \quad \text{and} \quad \|d_{k+1}\| \leq \kappa\|g_{k+1}\|^q.
\]

From the entry condition (Section~\ref{subsubsec:entry}), for all $k \geq k_0$ such that $x_k \in B(x^\star, \delta)$:
\[
\|g_k\| \leq L_1\delta.
\]

The step length satisfies:
\[
\|x_{k+1} - x_k\| = \alpha_k\|d_k\| \leq \alpha_k \kappa\|g_k\|^q.
\]
Substituting the optimal stepsize $\alpha_k^\star = \frac{\sigma}{L_1\kappa^2}\|g_k\|^{1+p-2q}$:
\begin{align*}
\|x_{k+1} - x_k\| &\leq \frac{\sigma}{L_1\kappa^2}\|g_k\|^{1+p-2q} \cdot \kappa\|g_k\|^q \\
&= \frac{\sigma}{L_1\kappa}\|g_k\|^{1+p-q}.
\end{align*}
Using $\|g_k\| \leq L_1\delta$:
\begin{equation}
\label{eq:step_NR}
\|x_{k+1} - x_k\| \leq \frac{\sigma}{L_1\kappa}(L_1\delta)^{1+p-q} = \frac{\sigma}{\kappa}L_1^{p-q}\,\delta^{1+p-q}.
\end{equation}

We want $\|x_{k+1} - x_k\| \leq \delta$, i.e.:
\[
\frac{\sigma}{\kappa}L_1^{p-q}\,\delta^{1+p-q} \leq \delta \iff \frac{\sigma}{\kappa}(L_1\delta)^{p-q} \leq 1.
\]

Three cases arise depending on the sign of $p - q$:

\paragraph{Case $p = q$:} The condition becomes $\frac{\sigma}{\kappa} \leq 1$, which always holds since $\sigma \leq 1 \leq \kappa$.

\paragraph{Case $p > q$ ($p - q > 0$):} The condition becomes:
\[
(L_1\delta)^{p-q} \leq \frac{\kappa}{\sigma} \iff \delta \leq \left(\frac{\kappa}{\sigma}\right)^{1/(p-q)}\frac{1}{L_1}.
\]
This is a condition for $\delta$ small enough, which is satisfied for sufficiently small $\delta$.

\paragraph{Case $p < q$ ($p - q < 0$):} Since $p - q < 0$, we have $(L_1\delta)^{p-q} = \frac{1}{(L_1\delta)^{q-p}} \to \infty$ as $\delta \to 0$. The condition becomes:
\[
(L_1\delta)^{p-q} \leq \frac{\kappa}{\sigma} \iff \delta^{p-q} \leq \frac{\kappa}{\sigma}L_1^{q-p} \iff \delta \geq \left(\frac{\sigma}{\kappa}\right)^{1/(q-p)}\frac{1}{L_1}.
\]
This is a condition for $\delta$ \emph{large enough}: the neighborhood cannot be too small. This creates a tension with the finite termination result (Case~3 in Section~\ref{subsec:regime3}), which also requires $p < q$.

\subsubsection{Restarted case ($k \in \mathcal{R}$)}

For restarted iterations, $d_k = -g_k$ and $\alpha_k = 1/L_1$. The step length is:
\[
\|x_{k+1} - x_k\| = \frac{1}{L_1}\|g_k\| \leq \frac{1}{L_1} \cdot L_1\delta = \delta.
\]
No condition on any parameter is needed: the step length is always bounded by $\delta$.

By the triangle inequality, for both cases:
\[
\|x_{k+1} - x^\star\| \leq \|x_{k+1} - x_k\| + \|x_k - x^\star\| \leq \delta + \delta = 2\delta.
\]
Therefore $x_{k+1} \in B(x^\star, 2\delta)$, completing the induction.
\end{proof}

\subsection{Are we staying in Regime 3 (Same proof technique as in \ref{subsubsec:staying})}
By using the exact same proof technique as in \ref{subsubsec:staying}, but replacing $B(x^\star, 2\delta)$ by $B(x^\star, \delta)$, we find the following condition in the contradiction:
\begin{equation}
    L_1 \geq \gamma .
\end{equation}

This is always true, meaning that we can prove that we stay in Regime~3 only for $L_1 = \gamma$, which is the most well-posed a problem can be.
This proof technique only allows to prove that we stay in Regime 3 once we enter it for perfect quadratic (i.e. condition number is 1) functions, that is 
function of the form $f(x) = x^T A x + b^T x +c$ with $A = \gamma I$. This is a well-known not interesting result.


\subsection{Staying in Regime 3 (Using gradient norm bounds)}

We attempt to prove that iterates eventually enter and stay in $B(x^\star, \delta)$, not just $B(x^\star, 2\delta)$, by showing that gradient norms shrink, which in turn shrinks step lengths and distances to $x^\star$.

\subsubsection{Restarted case ($k \in \mathcal{R}$)}

For restarted iterations, $d_k = -g_k$, $\alpha_k = 1/L_1$, so $x_{k+1} - x_k = -g_k/L_1$. We expand the gradient at $x_{k+1}$ using a Taylor expansion:
\[
\nabla f(x_{k+1}) = \nabla f(x_k) + \nabla^2 f(x_k)(x_{k+1} - x_k) + R_k,
\]
where $R_k$ is the remainder. Substituting $x_{k+1} - x_k = -g_k/L_1$:
\[
g_{k+1} = g_k - \frac{1}{L_1}\nabla^2 f(x_k)g_k + R_k = \left(I - \frac{\nabla^2 f(x_k)}{L_1}\right)g_k + R_k.
\]

\paragraph{Bounding the remainder.} We write:
\[
R_k = g_{k+1} - g_k - \nabla^2 f(x_k)(x_{k+1} - x_k) = \int_0^1 \left[\nabla^2 f(x_k + t(x_{k+1} - x_k)) - \nabla^2 f(x_k)\right](x_{k+1} - x_k)\, dt.
\]
Taking norms and using A2 ($\nabla^2 f$ is $L_2$-Lipschitz):
\[
\|R_k\| \leq \int_0^1 \|\nabla^2 f(x_k + t(x_{k+1} - x_k)) - \nabla^2 f(x_k)\| \cdot \|x_{k+1} - x_k\|\, dt \leq \int_0^1 L_2\, t\, \|x_{k+1} - x_k\|^2\, dt = \frac{L_2}{2}\|x_{k+1} - x_k\|^2.
\]
Since $\|x_{k+1} - x_k\| = \|g_k\|/L_1$:
\[
\|R_k\| \leq \frac{L_2}{2L_1^2}\|g_k\|^2.
\]

\paragraph{Gradient contraction.} Taking norms in the Taylor expansion:
\[
\|g_{k+1}\| \leq \left\|I - \frac{\nabla^2 f(x_k)}{L_1}\right\|\|g_k\| + \|R_k\|.
\]
In the strongly convex region $B(x^\star, 2\delta)$, the eigenvalues of $\nabla^2 f(x_k)$ lie in $[\gamma, L_1]$, so the eigenvalues of $I - \nabla^2 f(x_k)/L_1$ lie in $[0, 1 - \gamma/L_1]$. Therefore:
\begin{equation}
\label{eq:grad_contraction}
\|g_{k+1}\| \leq \left(1 - \frac{\gamma}{L_1}\right)\|g_k\| + \frac{L_2}{2L_1^2}\|g_k\|^2.
\end{equation}

\paragraph{Condition for gradient shrinkage.} We need the right-hand side of~\eqref{eq:grad_contraction} to be strictly less than $\|g_k\|$:
\[
\left(1 - \frac{\gamma}{L_1}\right)\|g_k\| + \frac{L_2}{2L_1^2}\|g_k\|^2 < \|g_k\| \iff \frac{L_2}{2L_1^2}\|g_k\| < \frac{\gamma}{L_1}.
\]
In the strongly convex region, $\|g_k\| \leq L_1\delta$ (from Section~\ref{subsubsec:entry}). Substituting:
\[
\frac{L_2}{2L_1^2} \cdot L_1\delta < \frac{\gamma}{L_1} \iff \frac{L_2}{2L_1}\delta < \frac{\gamma}{L_1} \iff \delta < \frac{2\gamma}{L_2}.
\]
Under this condition, gradients shrink at every iteration.

\paragraph{Step length and distance to $x^\star$.} Since gradients shrink, step lengths also shrink:
\[
\|x_{k+1} - x_k\| = \frac{\|g_k\|}{L_1} \leq \delta.
\]
From strong monotonicity ($\|g_k\| \geq \gamma\|x_k - x^\star\|$, see Section~\ref{subsubsec:entry}):
\[
\|x_k - x^\star\| \leq \frac{\|g_k\|}{\gamma}.
\]
By the triangle inequality:
\[
\|x_{k+1} - x^\star\| \leq \|x_{k+1} - x_k\| + \|x_k - x^\star\| \leq \frac{\|g_k\|}{L_1} + \frac{\|g_k\|}{\gamma} = \|g_k\|\left(\frac{1}{L_1} + \frac{1}{\gamma}\right).
\]
For this to be $\leq \delta$, we would need:
\[
L_1\delta\left(\frac{1}{L_1} + \frac{1}{\gamma}\right) \leq \delta \iff 1 + \frac{L_1}{\gamma} \leq 1 \iff \frac{L_1}{\gamma} \leq 0,
\]
which is impossible. Therefore, the first step after entering $B(x^\star, \delta)$ may leave this ball.

\paragraph{Eventual convergence to $B(x^\star, \delta)$.} Although the first step may leave $B(x^\star, \delta)$, the iterate remains in $B(x^\star, 2\delta)$ (from Proposition~\ref{prop:stay} or~\ref{prop:stay_step}), so the gradient contraction~\eqref{eq:grad_contraction} continues to apply. Under $\delta < 2\gamma/L_2$, the effective contraction rate is:
\[
\bar{\rho} := 1 - \frac{\gamma}{L_1} + \frac{L_2\delta}{2L_1} < 1.
\]
Iterating from $k_0$:
\[
\|g_{k_0+m}\| \leq \bar{\rho}^m \|g_{k_0}\| \leq \bar{\rho}^m L_1\delta.
\]
The distance to $x^\star$ satisfies:
\[
\|x_{k_0+m} - x^\star\| \leq \frac{\|g_{k_0+m}\|}{\gamma} \leq \frac{\bar{\rho}^m L_1\delta}{\gamma}.
\]
The step length at iteration $m$:
\[
\|x_{k_0+m+1} - x_{k_0+m}\| \leq \frac{\|g_{k_0+m}\|}{L_1} \leq \bar{\rho}^m\delta.
\]
By the triangle inequality:
\[
\|x_{k_0+m+1} - x^\star\| \leq \bar{\rho}^m\delta + \frac{\bar{\rho}^m L_1\delta}{\gamma} = \bar{\rho}^m\delta\left(1 + \frac{L_1}{\gamma}\right).
\]
This is $\leq \delta$ when:
\[
\bar{\rho}^m\left(1 + \frac{L_1}{\gamma}\right) \leq 1.
\]
Taking logarithms and using $-\log\bar{\rho} \geq \frac{2\gamma - L_2\delta}{2L_1}$:
\begin{equation}
\label{eq:transient}
m^\star_{\mathcal{R}} = \left\lceil\frac{2L_1\log\left(1 + L_1/\gamma\right)}{2\gamma - L_2\delta}\right\rceil.
\end{equation}
After $m^\star_{\mathcal{R}}$ restarted iterations in the strongly convex region, all subsequent iterates remain in $B(x^\star, \delta)$.

\paragraph{Summary.} Under $\delta < 2\gamma/L_2$:
\begin{enumerate}
    \item Gradient norms shrink geometrically: $\|g_{k+1}\| \leq \bar{\rho}\|g_k\|$ with $\bar{\rho} < 1$.
    \item Iterates stay in $B(x^\star, 2\delta)$ from the first iteration $k_0$.
    \item After a transient phase of $m^\star_{\mathcal{R}}$ iterations~\eqref{eq:transient}, iterates enter and stay in $B(x^\star, \delta)$.
\end{enumerate}

\subsubsection{Non-Restarted Case ($k \in \mathcal{N}$)}

We present two attempts to prove gradient shrinkage for non-restarted iterations. Both fail, for instructive reasons.

\paragraph{Attempt 1: Taylor expansion approach.} Recall that for non-restarted iterations, $d_k = -g_k + \beta_k d_{k-1}$ and $\alpha_k^\star = \frac{\sigma}{L_1\kappa^2}\|g_k\|^{1+p-2q}$. The step is $x_{k+1} - x_k = \alpha_k d_k$, so $x_{k+1} - x_k = \alpha_{k+1}d_{k+1}$ at the next iteration. Expanding the gradient at $x_{k+1}$ by Taylor:
\[
g_{k+1} = g_k + \nabla^2 f(x_k)(x_{k+1} - x_k) + R_k = g_k + \alpha_k\nabla^2 f(x_k) d_k + R_k.
\]
Substituting $\alpha_k^\star$:
\[
g_{k+1} = g_k + \frac{\sigma}{L_1\kappa^2}\|g_k\|^{1+p-2q}\nabla^2 f(x_k) d_k + R_k.
\]
Using $\|g_k\| \leq L_1\delta$:
\[
g_{k+1} \leq g_k + \frac{\sigma}{\kappa^2}(L_1\delta)^{p-2q}\nabla^2 f(x_k) d_k + R_k.
\]
This does not simplify further. In the restarted case, $d_k = -g_k$ allowed the factorization $g_k + \alpha_k\nabla^2 f(x_k)d_k = (I - \frac{\nabla^2 f(x_k)}{L_1})g_k$, yielding eigenvalue-based contraction. Here, $d_k = -g_k + \beta_k d_{k-1}$ introduces the term $\nabla^2 f(x_k)d_{k-1}$, which couples the current Hessian with the previous direction. The eigenvalue argument does not apply because $d_{k-1}$ is not aligned with any eigenvector of $\nabla^2 f(x_k)$ in general.

\paragraph{Attempt 2: Function value approach.} We use the smoothness bound $\|g_{k+1}\|^2 \leq 2L_1(f(x_{k+1}) - f(x^\star))$ and the decrease recursion~\eqref{eq:contraction_N3}:
\[
\|g_{k+1}\|^2 \leq 2L_1(f(x_{k+1}) - f(x^\star)) \leq 2L_1(f(x_k) - f(x^\star))\left(1 - C(f(x_k) - f(x^\star))^{r-1}\right).
\]
From PL, $f(x_k) - f(x^\star) \leq \frac{\|g_k\|^2}{2\gamma}$. Substituting:
\[
\|g_{k+1}\|^2 \leq 2L_1 \cdot \frac{\|g_k\|^2}{2\gamma}\left(1 - C\left(\frac{\|g_k\|^2}{2\gamma}\right)^{r-1}\right) = \frac{L_1}{\gamma}\|g_k\|^2\left(1 - C\left(\frac{\|g_k\|^2}{2\gamma}\right)^{r-1}\right).
\]
Substituting $C = \frac{\sigma^2(2\gamma)^r}{2L_1\kappa^2}$:
\[
\|g_{k+1}\|^2 \leq \frac{L_1}{\gamma}\|g_k\|^2\left(1 - \frac{\sigma^2(2\gamma)^r}{2L_1\kappa^2}\cdot\frac{\|g_k\|^{2(r-1)}}{(2\gamma)^{r-1}}\right) = \frac{L_1}{\gamma}\|g_k\|^2\left(1 - \frac{\sigma^2\gamma}{L_1\kappa^2}\|g_k\|^{2(r-1)}\right).
\]
The contraction factor for $\|g_{k+1}\|^2$ is therefore:
\[
\rho_k = \frac{L_1}{\gamma} - \frac{\sigma^2}{\kappa^2}\|g_k\|^{2(r-1)}.
\]
For $r = 1$ ($p = q$), this gives $\rho_k = \frac{L_1}{\gamma} - \frac{\sigma^2}{\kappa^2}$. Gradient shrinkage requires $\rho_k < 1$, i.e.:
\[
\frac{\sigma^2}{\kappa^2} > \frac{L_1 - \gamma}{\gamma} = \kappa_f - 1.
\]
Since $\sigma \leq 1$ and $\kappa \geq 1$, we have $\frac{\sigma^2}{\kappa^2} \leq 1$. But $\kappa_f - 1 > 1$ whenever $\kappa_f > 2$, which is the case for any moderately ill-conditioned problem. The condition is therefore not satisfiable in general.

\subsubsection{Non-Restarted Case ($k \in \mathcal{N}$) — Defining $d_k = -H_k g_k$}

We write the non-restarted direction as $d_k = -H_k g_k$, where $H_k$ is a positive definite matrix defined implicitly by the conjugate gradient update $d_k = -g_k + \beta_k d_{k-1}$. We denote the eigenvalues of $H_k$ by $\lambda_{\min}^H$ and $\lambda_{\max}^H$.

\paragraph{Step 1: Taylor expansion.} Expanding the gradient at $x_{k+1}$:
\[
g_{k+1} = g_k + \nabla^2 f(x_k)(x_{k+1} - x_k) + R_k.
\]
Since $x_{k+1} - x_k = \alpha_k d_k = -\alpha_k H_k g_k$:
\[
g_{k+1} = g_k - \alpha_k \nabla^2 f(x_k) H_k g_k + R_k = \left(I - \alpha_k H_k \nabla^2 f(x_k)\right)g_k + R_k.
\]

\paragraph{Step 2: Bound the remainder.} Using A2 and the fundamental theorem of calculus:
\[
\|R_k\| \leq \frac{L_2}{2}\|x_{k+1} - x_k\|^2 = \frac{L_2}{2}\|\alpha_k H_k g_k\|^2 \leq \frac{L_2\alpha_k^2(\lambda_{\max}^H)^2}{2}\|g_k\|^2.
\]
Substituting $\alpha_k = \frac{\sigma}{L_1\kappa^2}\|g_k\|^{1+p-2q}$:
\begin{equation}
\label{eq:Rk_NR}
\|R_k\| \leq \frac{L_2(\sigma\lambda_{\max}^H)^2}{2L_1^2\kappa^4}\|g_k\|^{2(2+p-2q)}.
\end{equation}

\paragraph{Step 3: Gradient contraction.} Taking norms in the Taylor expansion:
\[
\|g_{k+1}\| \leq \left\|I - \alpha_k H_k \nabla^2 f(x_k)\right\|\|g_k\| + \|R_k\|.
\]
In $B(x^\star, 2\delta)$, eigenvalues of $\nabla^2 f(x_k)$ lie in $[\gamma, L_1]$ and those of $H_k$ in $[\lambda_{\min}^H, \lambda_{\max}^H]$. The eigenvalues of $I - \alpha_k H_k \nabla^2 f(x_k)$ lie in:
\[
\left[1 - \alpha_k\lambda_{\max}^H L_1,\; 1 - \alpha_k\lambda_{\min}^H\gamma\right].
\]
Assuming all eigenvalues are nonnegative ($\alpha_k\lambda_{\max}^H L_1 \leq 1$), the operator norm is:
\[
\left\|I - \alpha_k H_k \nabla^2 f(x_k)\right\| = 1 - \alpha_k\lambda_{\min}^H\gamma = 1 - \frac{\sigma\lambda_{\min}^H\gamma}{L_1\kappa^2}\|g_k\|^{1+p-2q}.
\]
Combining with~\eqref{eq:Rk_NR}:
\begin{equation}
\label{eq:grad_contraction_NR_full}
\|g_{k+1}\| \leq \left(1 - \frac{\sigma\lambda_{\min}^H\gamma}{L_1\kappa^2}\|g_k\|^{1+p-2q}\right)\|g_k\| + \frac{L_2(\sigma\lambda_{\max}^H)^2}{2L_1^2\kappa^4}\|g_k\|^{2(2+p-2q)}.
\end{equation}
Factoring out $\|g_k\|$:
\[
\|g_{k+1}\| \leq \|g_k\|\left(1 - \frac{\sigma\lambda_{\min}^H\gamma}{L_1\kappa^2}\|g_k\|^{1+p-2q} + \frac{L_2(\sigma\lambda_{\max}^H)^2}{2L_1^2\kappa^4}\|g_k\|^{2(2+p-2q)-1}\right).
\]
Letting $u := \|g_k\|^{1+p-2q}$, we note that $2(2+p-2q)-1 = (1+p-2q) + (2+p-2q)$, so $\|g_k\|^{2(2+p-2q)-1} = u \cdot \|g_k\|^{2+p-2q}$. The contraction factor becomes:
\begin{equation}
\label{eq:rho_k_NR}
\rho_k := 1 - u\left(\frac{\sigma\lambda_{\min}^H\gamma}{L_1\kappa^2} - \frac{L_2(\sigma\lambda_{\max}^H)^2}{2L_1^2\kappa^4}\|g_k\|^{2+p-2q}\right).
\end{equation}

\paragraph{Step 4: Condition for gradient shrinkage.} We need $\rho_k < 1$, i.e., the term in parentheses must be positive:
\[
\frac{\sigma\lambda_{\min}^H\gamma}{L_1\kappa^2} > \frac{L_2(\sigma\lambda_{\max}^H)^2}{2L_1^2\kappa^4}\|g_k\|^{2+p-2q}.
\]
Rearranging:
\[
\|g_k\|^{2+p-2q} < \frac{2\sigma\lambda_{\min}^H\gamma L_1}{\kappa^2 L_2(\lambda_{\max}^H)^2}.
\]
This must hold for all $k \geq k_0$ with $x_k \in B(x^\star, 2\delta)$. At entry, $\|g_{k_0}\| \leq L_1\delta$, and since we will show gradients shrink, the worst case is at entry. So it suffices to require:
\begin{equation}
\label{eq:delta_cond_NR_tight}
(L_1\delta)^{2+p-2q} < \frac{2\sigma\lambda_{\min}^H\gamma L_1}{\kappa^2 L_2(\lambda_{\max}^H)^2} \iff \delta < \frac{1}{L_1}\left(\frac{2\sigma\lambda_{\min}^H\gamma L_1}{\kappa^2 L_2(\lambda_{\max}^H)^2}\right)^{\frac{1}{2+p-2q}}.
\end{equation}

\paragraph{Step 5: Effective contraction rate.} Under condition~\eqref{eq:delta_cond_NR_tight}, $\rho_k < 1$ for all $k \geq k_0$. Moreover, since $\|g_{k+1}\| \leq \rho_k\|g_k\| < \|g_k\|$, the contraction rate $\rho_k$ \emph{improves} at each step (since smaller $\|g_k\|$ makes the negative terms in~\eqref{eq:rho_k_NR} smaller, but the positive term $1$ stays). Therefore, the worst-case rate is at the first iteration:
\begin{equation}
\label{eq:rho_bar_NR}
\bar{\rho} := 1 - (L_1\delta)^{1+p-2q}\left(\frac{\sigma\lambda_{\min}^H\gamma}{L_1\kappa^2} - \frac{L_2(\sigma\lambda_{\max}^H)^2}{2L_1^2\kappa^4}(L_1\delta)^{2+p-2q}\right) < 1.
\end{equation}
By induction, for all $m \geq 0$:
\begin{equation}
\label{eq:grad_decay_NR_tight}
\|g_{k_0+m}\| \leq \bar{\rho}^m\|g_{k_0}\| \leq \bar{\rho}^m L_1\delta.
\end{equation}

\paragraph{Step 6: Step length after $m$ iterations.} Using $\|d_k\| = \|H_k g_k\| \leq \lambda_{\max}^H\|g_k\|$ and $\alpha_k = \frac{\sigma}{L_1\kappa^2}\|g_k\|^{1+p-2q}$:
\[
\|x_{k_0+m+1} - x_{k_0+m}\| = \alpha_k\|d_k\| \leq \frac{\sigma\lambda_{\max}^H}{L_1\kappa^2}\|g_{k_0+m}\|^{2+p-2q}.
\]
Substituting~\eqref{eq:grad_decay_NR_tight}:
\begin{equation}
\label{eq:step_NR_tight}
\|x_{k_0+m+1} - x_{k_0+m}\| \leq \frac{\sigma\lambda_{\max}^H}{L_1\kappa^2}(\bar{\rho}^m L_1\delta)^{2+p-2q} = \frac{\sigma\lambda_{\max}^H(L_1\delta)^{2+p-2q}}{L_1\kappa^2}\bar{\rho}^{m(2+p-2q)}.
\end{equation}
Since $2+p-2q > 1$, this decays as $\bar{\rho}^{m(2+p-2q)}$, which is faster than $\bar{\rho}^m$.

\paragraph{Step 7: Distance to $x^\star$ after $m$ iterations.} From strong monotonicity ($\|g_k\| \geq \gamma\|x_k - x^\star\|$):
\begin{equation}
\label{eq:dist_NR_tight}
\|x_{k_0+m} - x^\star\| \leq \frac{\|g_{k_0+m}\|}{\gamma} \leq \frac{\bar{\rho}^m L_1\delta}{\gamma}.
\end{equation}

\paragraph{Step 8: Triangle inequality.} Combining Steps 6 and 7:
\begin{align*}
\|x_{k_0+m+1} - x^\star\| &\leq \|x_{k_0+m+1} - x_{k_0+m}\| + \|x_{k_0+m} - x^\star\| \\
&\leq \frac{\sigma\lambda_{\max}^H(L_1\delta)^{2+p-2q}}{L_1\kappa^2}\bar{\rho}^{m(2+p-2q)} + \frac{\bar{\rho}^m L_1\delta}{\gamma}.
\end{align*}
Since $2+p-2q > 1$, for any $m \geq 1$ we have $\bar{\rho}^{m(2+p-2q)} < \bar{\rho}^m$, so the first term is dominated by the second. Bounding:
\[
\|x_{k_0+m+1} - x^\star\| \leq \frac{\bar{\rho}^m L_1\delta}{\gamma}\left(1 + \frac{\sigma\lambda_{\max}^H\gamma(L_1\delta)^{1+p-2q}}{L_1\kappa^2}\right).
\]
Define:
\begin{equation}
\label{eq:A_NR}
A := \frac{L_1}{\gamma}\left(1 + \frac{\sigma\lambda_{\max}^H\gamma}{L_1\kappa^2}(L_1\delta)^{1+p-2q}\right).
\end{equation}

\paragraph{Step 9: Solve for $m$.} We need $\|x_{k_0+m+1} - x^\star\| \leq \delta$, i.e.:
\[
\bar{\rho}^m \cdot A \leq 1.
\]
Taking logarithms. Since $\bar{\rho} < 1$, $\log\bar{\rho} < 0$, dividing flips the inequality:
\[
m \geq \frac{\log A}{-\log\bar{\rho}}.
\]
To bound $-\log\bar{\rho}$, write $\bar{\rho} = 1 - X$ where:
\[
X := (L_1\delta)^{1+p-2q}\left(\frac{\sigma\lambda_{\min}^H\gamma}{L_1\kappa^2} - \frac{L_2(\sigma\lambda_{\max}^H)^2}{2L_1^2\kappa^4}(L_1\delta)^{2+p-2q}\right) > 0.
\]
Using $-\log(1 - X) \geq X$ for $X \in (0,1)$:
\begin{equation}
\label{eq:mstar_NR_tight}
\boxed{m^\star_{\mathcal{N}} \leq \left\lceil\frac{\log A}{(L_1\delta)^{1+p-2q}\left(\dfrac{\sigma\lambda_{\min}^H\gamma}{L_1\kappa^2} - \dfrac{L_2(\sigma\lambda_{\max}^H)^2}{2L_1^2\kappa^4}(L_1\delta)^{2+p-2q}\right)}\right\rceil}
\end{equation}
where $A$ is defined in~\eqref{eq:A_NR}, under the condition~\eqref{eq:delta_cond_NR_tight} on $\delta$.

\paragraph{Verification: recovery of the restarted case.} For $H_k = I$ ($\lambda_{\min}^H = \lambda_{\max}^H = 1$) and $p = q = 1$ ($1+p-2q = 0$, $2+p-2q = 1$):
\begin{itemize}
    \item $(L_1\delta)^{1+p-2q} = 1$ and $(L_1\delta)^{2+p-2q} = L_1\delta$.
    \item The denominator becomes $\frac{\sigma\gamma}{L_1\kappa^2} - \frac{L_2\sigma^2}{2L_1^2\kappa^4}L_1\delta$. For $\sigma = \kappa = 1$: $\frac{\gamma}{L_1} - \frac{L_2\delta}{2L_1} = \frac{2\gamma - L_2\delta}{2L_1}$.
    \item $A = \frac{L_1}{\gamma}(1 + \frac{\sigma\gamma}{L_1\kappa^2}) \approx \frac{L_1}{\gamma}$ for typical parameters, so $\log A \approx \log(L_1/\gamma)$.
    \item This gives $m^\star \approx \frac{2L_1\log(L_1/\gamma)}{2\gamma - L_2\delta}$, matching $m^\star_{\mathcal{R}}$ from~\eqref{eq:transient}.
\end{itemize}

\begin{remark}
The eigenvalues $\lambda_{\min}^H$ and $\lambda_{\max}^H$ of the implicit preconditioner $H_k$ are not known in general for nonlinear conjugate gradient. Bounding them requires assumptions on the conjugate direction parameter $\beta_k$ and is specific to each CG variant (FR, PRP+, HZ, etc.). For a specific $\beta_k$ formula, these eigenvalues can be estimated, yielding a concrete expression for $m^\star_{\mathcal{N}}$.
\end{remark}

\bibliographystyle{plain}
\bibliography{refs} 

\end{document}